\section{DISCUSSION}

\begin{enumerate}
    \item \textbf{Pressure Drop Trends:}
    Based on the experimental data and charts:
    \begin{itemize}
        \item \textbf{Dry Column ($L=0$):} As gas flow rate $G$ increases, the pressure drop $\Delta P$ increases approximately linearly on a log-log scale, following the theoretical prediction $\Delta P \propto G^n$.
        \item \textbf{Wetted Column ($L>0$):} As $G$ increases, $\Delta P$ increases more rapidly than in the dry case. As the liquid irrigation rate $L$ increases, the column reaches the flooding point at lower gas velocities. This is because the liquid occupies the void space, increasing the actual gas velocity and resistance.
    \end{itemize}

    \item \textbf{Friction Factor and Reynolds Number:}
    The friction factor $f$ was found to decrease as the Reynolds number $Re$ increases, which is typical for turbulent flow in packed beds. The chart of $f$ vs. $Re$ helps determine the optimal operating range to maximize mass transfer efficiency while maintaining a manageable pressure drop.

    \item \textbf{Comparison with Predictions:}
    The relationship between $\log(\Delta P/Z)$ and $\log G$ is linear as predicted. For wetted columns, two distinct regions are observed: below and above the loading point. Below the loading point, $\Delta P$ increases gradually. Above the loading point, $\Delta P$ rises sharply as the column approaches flooding.

    \item \textbf{Sources of Error:}
    Several factors could contribute to discrepancies between experimental results and theoretical predictions:
    \begin{itemize}
        \item \textbf{Flow Instability}: The liquid pump and the air blower (fan) may produce unstable flow rates, leading to fluctuating pressure readings.
        \item \textbf{Bottom Water Level}: If the water column maintained at the bottom of the column is insufficient, it can lead to flooding in the differential pressure tube, affecting the results.
        \item \textbf{Measurement Errors}: Human error in reading the manometer scales and flowmeters, especially under fluctuating conditions.
        \item \textbf{Equilibrium Time}: Not allowing sufficient time (suggested 10 minutes) for the system to stabilize after adjusting flow rates.
        \item \textbf{Ambient Conditions}: Variations in temperature and pressure during the experiment affecting air density and viscosity.
    \end{itemize}
\end{enumerate}
